\subsection{An independently sized complex network}
\label{sec:independent-complex}
Keep the paired bit network at $h_{\rm b}=50$, with address-swap exponent
$\tau=1-296/10^{11}$. Instantiate the original complex construction instead
at $h_{\rm c}=25$. Its constants, distinguished from those of the bit
network, are
\[
 m_{\rm c}=15625,\quad v_{\rm c}=2300,\quad N_{\rm c}=12167000000,
\]
\[
 W_{\rm c}=58645352620000,\quad L_{\rm c}=10315500000,\quad
 s_{\rm c}=916333630984500000.
\]
Here $W_{\rm c}=2v_{\rm c}^3+3v_{\rm c}^2(v_{\rm c}z_{\rm c}+26)$,
$z_{\rm c}=\binom{22}{3}+66$, and $L_{\rm c}=3v_{\rm c}^2\cdot25\cdot26$.

\begin{proposition}\label{prop:compact-complex-interface}
The original complex phase interface, including its signed scalar bank
exchange, arbitrary-scratch restoration, binary residual bases and endpoint
corrections, holds with these constants. Its normalized residual saving
supports the strict exponent $\sigma=1-418/10^{12}$.
\end{proposition}
\begin{proof}
The original construction and residual table apply at this ground size.
For an intersection of size $j$, the center coefficient is $(j-1)/2$; the
side correction is its negative for distinct even-intersection neighbors.
Their sum is zero for $j=0,1,2$ and one for $j=3$. The same eight-step
schedule therefore restores arbitrary scratch and gives the elementary bank
shear, and the three stages give the signed exchange.

Binary triple indicators have norm one. The orthogonal decompositions in
the residual table remain nondegenerate. Every nonzero residual has a
norm-one vector: choose a coordinate outside the support of a triple or two
neighboring triples (at most six coordinates), or outside an earlier/future
tensor indicator (support $3^j<25^j$); tensor with the remaining norm-one
factors. Thus the original nonalternating-space argument supplies the required
orthonormal bases. No even-ground-size matching or rational bit-interface
positivity is needed for this unchanged complex network.

The rank accounting gives
\[
 s_{\rm c}=W_{\rm c}m_{\rm c}-2N_{\rm c}+2L_{\rm c},\qquad
 \eta_{\rm c}=\frac{W_{\rm c}m_{\rm c}-s_{\rm c}}{W_{\rm c}m_{\rm c}}
             =\frac{14}{3464399375}>0.
\]
The needed positivity is $L_{\rm c}<N_{\rm c}$, not
$2L_{\rm c}<N_{\rm c}$. Endpoint weight-27 phase corrections and the
signed-exchange correction are unchanged. Exact rational logarithm bounds
give $\log m_{\rm c}<966/100$ and
$\eta_{\rm c}>(418/10^{12})(966/100)$. The inequality
$e^{-x}>1-x$ then yields $s_{\rm c}/W_{\rm c}<m_{\rm c}^{\sigma}$.
The certificate also checks $2\le s_{\rm c}<m_{\rm c}^5$.
\end{proof}

The bit and complex dimensions need not agree. The bit interface supplies
swaps at arbitrary binary widths. The complex recursion uses only that
interface's exponent $\tau$ and fixed tape guarantee, while its own branching,
role split and depth use $m_{\rm c},W_{\rm c},s_{\rm c}$. All occurrences of
$m,W,s$ within the revised layer proof denote these complex constants.
The retained assembly consumes their time exponents and completed-layer
contract, so it introduces no equality requirement between the two arities.
